\section{Decoding Ablation}
\label{sec:ablation}

We fix the \crashlogic{} checkpoint and vary only the decoding rule.
Table~\ref{tab:e4} summarises the three configurations that matter for the claim; the full matrix (all $\alpha$, shuffle, SEASON) is Table~\ref{tab:e4-full}.

\begin{table*}[!t]
\centering
\caption{Decoding ablation on the fixed \crashlogic{} checkpoint ($n{=}150$). Mild reverse \tcd{} is safe but never significantly better than greedy; strong contrast significantly hurts (Table~\ref{tab:paired}). Shuffle and full SEASON rows are in Table~\ref{tab:e4-full}. $^{\dagger}$ marks a significant degradation vs.\ greedy ($p<0.05$).}
\label{tab:e4}
\small
\setlength{\tabcolsep}{10pt}
\begin{tabular*}{\textwidth}{@{\extracolsep{\fill}}lcccc@{}}
\toprule
\textbf{Variant} & \textbf{ROUGE-L}$\uparrow$ & \textbf{Full-NLI}$\uparrow$ & \textbf{\Hcost{}}$\downarrow$ & \textbf{\Ocost{}}$\downarrow$ \\
\midrule
Greedy ($\alpha{=}0$)           & 0.336 & \textbf{0.247} & 0.223 & 0.107 \\
\tcd{} $\alpha{=}0.5$ (reverse) & \textbf{0.339} & 0.234 & \textbf{0.212} & \textbf{0.102} \\
\tcd{} $\alpha{=}2.0$ (reverse) & 0.319$^{\dagger}$ & 0.148$^{\dagger}$ & 0.253$^{\dagger}$ & 0.128$^{\dagger}$ \\
\bottomrule
\end{tabular*}
\end{table*}


\paragraph{Contrast strength.}
From greedy to $\alpha{=}0.5$ every faithfulness metric moves in the desired direction by a small, non-significant amount.
At $\alpha{=}2.0$ hallucination, omission, ROUGE-L, and Full-NLI all degrade significantly (Table~\ref{tab:paired}), and the colour-disjoint agent error rises from $21\%$ to $30\%$ (Table~\ref{tab:fields}).
The trajectory matches prediction P3: beyond $\alpha{\approx}1$ the amplified component is mostly noise.

\paragraph{Shuffle and full SEASON.}
Neither alternative beats reverse \tcd{} on any faithfulness metric (Table~\ref{tab:e4-full}).
Shuffle posts the best lexical numbers and the worst NLI-Score; full SEASON raises both costs relative to greedy.
The failure to move \Hcost{} is therefore not a tuning problem: seven configurations spanning two negative constructions leave colour, weather, and severity conflict rates where greedy decoding put them.

\paragraph{What the ablation establishes.}
Mild temporal contrast is safe but ineffective; strong contrast is harmful; no recipe tested here reduces hallucination on this task.
